\newpage
%\thispagestyle{empty}
\begin{recipe}[source={La Mami},
	portion={4 personas},
	preparationtime={\unit[1]{Hora}}
	]{Fritos de atún}
%	\includegraphics[width=0.25\textwidth]{Fritos de atún}
	\introduction{
		Los fritos de atún son una receta típica de la familia. Son muy ricos para comer solos, o acompañándolos con algún arroz o ensalada.
	}
	\ingredients{
		\unit[500]{gr} & 500 gramos de lomitos de atún en agua, muy bien estilado, sin su jugo \\
		\unit[15]{gr} & de jengibre \\
		1 & ralladura de limón \\
		2 & cucharadas de perejil picado finito \\
		3 & cucharadas de la parte verde del cebollín picado \\
		1 o 2 & cucharadas de harina \\
		1 & cucharada de orégano fresco \\
		& Sal a gusto \\
		& Pimienta a gusto \\
		2 o 3 & huevos enteros, dependerá del tamaño de estos y de la humedad del atún
	}
	\preparation{
		\begin{enumerate}
			\item Comienza a realizar el batido, para ello, bate el huevo, mezcla la harina con el huevo y vierte agua poco a poco hasta conseguir una textura espesa, y reserva.
			\item En una olla o sartén, prepara un sofrito echando un poco del aceite con el ajo picándolo finamente y agregando la cebolla, el orégano y un poco de sal y revolviendo a fuego suave entre 5 y 15 minutos máximo.
			\item Mezcla el sofrito y la mezcla reservada previamente hasta conseguir una textura semilíquida.
			\item Mezclar el atún con el batido y preparar bolitas y luego aplastarlas un poco, deberían de quedar como de máximo \unit[1]{cm} de espesor.
			\item Unta en un poco de harina y pan rallado.
			\item En una sartén con aceite muy caliente, freír.
			\item Cuando estén doradas por ambos lados, retirar y dejar estilando en una fuente.
			\item Dejar enfriar unos minutos y servir.
		\end{enumerate}
	}
\end{recipe}